% How Roameo works inside, in three bands that match the three ideas. Six helpers and a
% coordinator. Four helpers use the AI model (backend/agents/intake.py, research.py, feedback.py,
% guide.py). The planner and the checker are plain code (composer.py, critic.py), and the planner
% uses the AI only to write the friendly summary. The coordinator is plain code (orchestrator.py)
% and stops after five fixes. In demo mode, simple rules stand in for the AI (backend/llm.py).
\documentclass[border=2pt]{standalone}
% Shared style of every Roameo figure (the deck and the poster).
% Type is near-black and grey on white. Teal, the colour of the app, is the brand and the system.
% One warm accent marks tension only (the budget gap, the over-budget plan). Green marks only a
% confirmed result: a human approval or a passed test. Thin rules, no shadows, no gradients.
\usepackage[T1]{fontenc}
\usepackage[scaled]{helvet}
\renewcommand{\familydefault}{\sfdefault}
\usepackage{tikz}
\usepackage{pgfplots}
\pgfplotsset{compat=1.18}
\usepackage{fontawesome5}
\hyphenpenalty=10000
\exhyphenpenalty=10000
\usetikzlibrary{arrows.meta,positioning,calc,fit,backgrounds,shapes.callouts,shapes.misc}

% Colours
\definecolor{primary}{HTML}{0F766E}   % Roameo teal: the system, the agents
\definecolor{primarysoft}{HTML}{E3F1EF}
\definecolor{accent}{HTML}{C2410C}    % tension only: the budget gap
\definecolor{accentsoft}{HTML}{FBE9E1}
\definecolor{ok}{HTML}{4D7C0F}        % a confirmed result: approval, a passed test
\definecolor{oksoft}{HTML}{EDF3E1}
\definecolor{ink}{HTML}{1F1F1F}
\definecolor{muted}{HTML}{5F6368}
\definecolor{quiet}{HTML}{9A9A9A}
\definecolor{rulegrey}{HTML}{C4C4C4}
\definecolor{gridgrey}{HTML}{E4E4E4}
\definecolor{faint}{HTML}{F4F4F4}

\tikzset{
  >={Stealth[length=2.4mm,width=1.8mm]},
  every picture/.style={line cap=butt, line join=miter},
  % text styles
  hd/.style={font=\bfseries, text=ink, align=left},
  sub/.style={text=muted, align=left},
  % shapes: 2 pt corners, thin rules
  card/.style={draw=rulegrey, line width=0.5pt, fill=white, rounded corners=2pt, inner sep=5pt, align=left, text=ink},
  sys/.style={draw=primary, line width=0.8pt, fill=white, rounded corners=2pt, inner sep=5pt, align=center, text=ink},
  sysfill/.style={draw=primary, line width=0.8pt, fill=primary, rounded corners=2pt, inner sep=5pt, align=center, text=white},
  code/.style={draw=ink, line width=0.8pt, fill=faint, rounded corners=2pt, inner sep=5pt, align=center, text=ink},
  warm/.style={draw=accent, line width=0.8pt, fill=accentsoft, rounded corners=2pt, inner sep=5pt, align=left, text=ink},
  receipt/.style={draw=ok, line width=0.8pt, fill=oksoft, rounded corners=2pt, inner sep=5pt, align=left, text=ink},
  planned/.style={draw=muted, line width=0.6pt, dash pattern=on 3pt off 2pt, fill=white, rounded corners=2pt, inner sep=5pt, align=left, text=muted},
  % chat bubbles: the traveller in teal, the app in light grey
  saybubble/.style={fill=primary, text=white, rounded corners=6pt, inner sep=5pt, align=left},
  appbubble/.style={fill=faint, text=ink, rounded corners=6pt, inner sep=5pt, align=left},
  chip/.style={draw=rulegrey, line width=0.5pt, fill=white, rounded corners=7pt, inner sep=4pt, text=ink},
  % connectors
  link/.style={->, draw=muted, line width=0.6pt},
  syslink/.style={->, draw=primary, line width=0.9pt},
  oklink/.style={->, draw=ok, line width=0.9pt},
  warmlink/.style={->, draw=accent, line width=0.9pt},
  nextlink/.style={->, draw=muted, line width=0.6pt, dash pattern=on 3pt off 2pt},
}

% Charts: a real axis, few light gridlines, no frame on the top and right.
\pgfplotsset{
  roameo/.style={
    axis line style={draw=muted, line width=0.5pt},
    tick style={draw=muted, line width=0.5pt},
    major tick length=3pt,
    tick label style={text=ink, /pgf/number format/assume math mode=true},
    label style={text=ink},
    title style={font=\bfseries, text=ink},
    grid style={draw=gridgrey, line width=0.4pt},
    every axis plot/.append style={line width=0.6pt},
    legend style={draw=none, fill=none, text=ink},
    legend cell align=left,
  },
}

\begin{document}
\begin{tikzpicture}[x=1cm, y=1cm, font=\small,
  helper/.style={fill=white, draw=umb, line width=0.8pt, rounded corners=3pt, inner sep=5pt, text width=2.3cm, align=flush left, text=ink, minimum height=1.6cm},
  person/.style={fill=white, draw=ok, line width=0.8pt, rounded corners=3pt, inner sep=5pt, text width=2.3cm, align=flush left, text=ink, minimum height=1.2cm},
  band/.style={rounded corners=4pt},
  bandlab/.style={anchor=north west, inner sep=0pt, font=\small\bfseries, text=umb, align=flush left, text width=2.2cm}]
  % bands
  \fill[umbsoft, band] (0,4.3) rectangle (17.2,6.55);
  \fill[faint, band] (0,2.0) rectangle (17.2,4.15);
  \fill[oksoft, band] (0,0) rectangle (17.2,1.85);
  \node[bandlab] at (0.2,6.35) {\markagents\enspace AI agents\\[1pt]\mdseries\footnotesize\textcolor{ink}{use the AI model, or simple rules in demo mode}};
  \node[bandlab] at (0.2,3.95) {\markcode\enspace Plain code\\[1pt]\mdseries\footnotesize\textcolor{ink}{the same sum every time}};
  \node[bandlab, text=ok] at (0.2,1.65) {\markyou\enspace You\\[1pt]\mdseries\footnotesize\textcolor{ink}{the final say}};
  % helpers
  \node[helper, anchor=north west] (reader) at (2.7,6.35) {\textbf{Request reader}\\ reads your sentence, asks what is missing};
  \node[helper, anchor=north west] (research) at (5.6,6.35) {\textbf{Researcher}\\ chooses the searches, at most 8 steps};
  \node[helper, anchor=north west] (change) at (11.4,6.35) {\textbf{Change reader}\\ turns your words into edits};
  \node[helper, anchor=north west] (guide) at (8.5,6.35) {\textbf{Guide}\\ writes a visit guide for each stop};
  \node[anchor=north west, text width=2.6cm, align=flush left, inner sep=0pt, text=ink, font=\footnotesize] at (14.4,6.25)
    {Four helpers use the AI model. Each has one narrow job, and code checks what it returns.};
  \node[helper, anchor=north west, draw=ink] (planner) at (5.6,3.95) {\textbf{Planner}\\ builds the days and prices every item};
  \node[helper, anchor=north west, draw=ink] (checker) at (8.5,3.95) {\textbf{Checker}\\ adds up the total, tries one fix at a time};
  \node[anchor=north west, text width=3.75cm, align=flush left, inner sep=0pt, text=ink, font=\footnotesize] at (13.35,3.95)
    {The \textbf{coordinator}, also plain code, passes the work along and stops after five fixes. The AI never adds up the money.};
  \node[person, anchor=north west] (type) at (2.7,1.7) {types a sentence, answers what is missing};
  \node[person, anchor=north west] (choose) at (8.5,1.7) {chooses how to close the gap};
  \node[person, anchor=north west] (ok) at (11.4,1.7) {asks for changes, or says OK};
  % flow
  \draw[syslink] (type.north) -- (reader.south);
  \draw[syslink] (reader.east) -- (research.west);
  \draw[syslink] (research.south) -- (planner.north);
  \draw[syslink] (planner.east) -- (checker.west);
  \draw[warmlink] ([xshift=-8pt]checker.south) -- ([xshift=-8pt]choose.north) node[midway, left, font=\footnotesize, text=accentink, inner sep=2pt] {still over};
  \draw[syslink] (choose.east) -- (ok.west);
  \draw[syslink] ([xshift=10pt]ok.north) -- ([xshift=10pt]change.south);
  \draw[syslink] ([xshift=-14pt]change.south) |- (checker.east) node[pos=0.25, left, font=\footnotesize, inner sep=2pt] {edit};
  \draw[oklink] (ok.east) -- ++(1.9,0) node[right, font=\footnotesize\bfseries, text=ok, align=flush left, inner sep=2pt, text width=1.7cm] {calendar file, after OK};
\end{tikzpicture}
\end{document}
